\subsection{Closing the fifth index}
\label{ext:five-assembly}
The fifth-index calculation has three complementary parts: a count
comparison at small third moments, an elementary corner estimate where
the signed moment is large and positive, and a Fourier estimate for the
remaining third moments. The corner is a genuine part of the cover;
it cannot be discarded as a small exceptional region.

\begin{proposition}[Fifth-index assembly criterion]
\label{ext:five-assembly-criterion}
Assume $C_4\le\CE$, the continuous inequalities proved in this paper,
and complete verification of the finite inputs and their implications
described below. Then $C_5\le\CE$.
\end{proposition}
\begin{proof}
Suppose $C_5>\CE$ and choose a positive closed-event ratio maximizer.
Its same-law predecessor score is at most $C_4$, so its contact
allowance is zero. The support calculations in
Appendix~\ref{ext:n5-support} therefore apply to this very marginal.
They precede the supported Fourier estimates. The unrestricted
auxiliary Fourier estimate first supplies $C_5<429/1000$; this is
the input to the two lowest-moment initializers, not their conclusion.
The higher-moment initializers use separately acquired supported
own-score bounds. Their scalar input ceilings are
\[
\begin{array}{c|cc}
\text{region}&\text{own score}&\text{HIGH}\\\hline
1\le\rho\le11/10&429/1000&337/1000\\
11/10\le\rho\le23/20&429/1000&369/1000\\
23/20\le\rho\le237/200&52/125&393/1000\\
237/200\le\rho\le6/5,\ \beta\le1/5&207/500&193/500\\
237/200\le\rho\le6/5,\ \beta\ge1/5&103/250&193/500
\end{array}
\]
Here HIGH is the compensated high-frequency contribution, as in
Lemma~\ref{ext:compensated-high}. Each such input is used only on a
support and moment domain contained in the domain of its Fourier
proof. In particular a better score bound requiring a smaller support
radius is not transferred to a larger-radius profile.
The last two rows are restricted to the count rectangles specified below;
the separate corner does not require these inputs.

Here is the exact signed-moment partition behind the small-moment
count calculation. Set
\[
 b_2=\frac{84729828287999}{250000000000000},\qquad
 b_3=\frac{35139006743001}{10^{14}}.
\]
The sharp signed-moment inequality bounds $|\beta|$ by $b_2$ on
$\rho\le237/200$ and by $b_3$ on $\rho\le6/5$.
The rational endpoints are checked by
$(1+b_i^2)^2\ge\rho_i^2(1-b_i^2)$ at the respective upper
third moments, using the monotonicity of this expression in $b_i\ge0$.

The two lowest moment intervals each form one full signed region.
On $[23/20,237/200]$, split $\rho$ at $467/400$ and $\beta$ at
$1/5$, giving four closed rectangles. On $[237/200,6/5]$, split
$\rho$ at $477/400$. In the lower half use the three signed intervals
$[-b_3,1/5]$, $[1/5,3/10]$, and $[3/10,b_3]$. In the upper half
use $[-b_3,1/5]$, $[1/5,3/10]$, and $[3/10,17/50]$.
These are the other six count regions. The remaining rectangle is
\begin{equation}\label{ext:five-corner-domain}
 [477/400,6/5]\times[17/50,b_3]
 \quad\text{in }(\rho,\beta).
\end{equation}
It is covered by the whole-line elementary estimate in
Appendix~\ref{ext:gapfree5}, including its two
unbounded threshold rays. This estimate has no support or predecessor
premise. Its upper bound is below $4097/10000<\CE$.

In each of the twelve count regions, an unseeded finite initializer
first acquires the raw remainder and then the physical defect, gap
and threshold constraints. The whole final moment cover is exported.
For each probability panel $[j/100,(j+1)/100]$, the subsequent resource
calculation concerns that same cover and that same law. Its signed
cubic, conditional variance and inward-displacement inequalities give
a joint count domain as in Section~\ref{ext:physical-count-domain}.
Every inequality in the numerical count contract must follow in its
stated direction from this acquired domain; matching a filename or
an isolated scalar maximum is insufficient.

The twelve regions have $1,200$ probability-panel instances. Of
these, $905$ are excluded by their resource inequalities. The other
$295$ are covered by six-root count forests with $1,770$ roots and
$75,178$ terminal leaves. Every accepted leaf has score strictly below
the fixed count target, and that target is strictly below $\CE$.
The interval boundaries in $\rho$, $\beta$ and $p$ are included.
Consequently these twelve regions and \eqref{ext:five-corner-domain}
cover all $1\le\rho\le6/5$, without assuming that the sign of
$\beta$ can be changed while keeping the event fixed.

For completeness, the supported Fourier complement has six closed
bands, each supplied with a previously acquired support radius:
\[
\begin{array}{c|c}
\rho& R\\\hline
{[6/5,5/4]}&13/5\\
{[5/4,7/5]}&117/40\\
{[7/5,3/2]}&25/8\\
{[3/2,8/5]}&133/40\\
{[8/5,17/10]}&7/2\\
{[17/10,7361/4000]}&15/4
\end{array}
\]
The first band uses the additional radius-$13/5$ support calculation.
The second restricts the wider support interval $[6/5,7/5]$.
All six implications retain both support signs and every real
threshold. The Fourier proof subdivides these bands into $41$ closed
strips and bounds the score below $4097/10000$ on each. Its recorded
$42$ mesh attempts include the unsuccessful coarser attempt; only
the successful complete enclosures are used as inequalities.
For $\rho\ge7361/4000$, the decreasing published affine bound is
at most $819/2000<4097/10000$.

Thus the count regions, the elementary corner, the six Fourier
bands and the infinite tail exclude the chosen maximizing witness.
Closed-event attainment and reflection account for either sign of
the original Kolmogorov discrepancy. This proves the implication.
\end{proof}

\subsection{Induction from six through fourteen}
